\chapter{Diseño y arquitectura de la solución}
\label{chap:diseno-arquitectura}

La arquitectura de Model2Blockly considera que un editor Blockly es una sintaxis
visual de un lenguaje de dominio específico. Por eso, el sistema debe partir de
una descripción abstracta del dominio y producir una implementación web ejecutable.

El capítulo se organiza de la siguiente manera. En primer lugar, se presenta el
enfoque arquitectónico y el flujo general de generación. A continuación, se
describe el modelo intermedio que conecta las entradas con los generadores.
Después se justifican las rutas basadas en Ecore y en el DSL textual
\texttt{.m2b}, así como la generación de los editores Blockly. Finalmente, se
exponen los mecanismos de validación, referencias y personalización, y se
relacionan las decisiones arquitectónicas con los objetivos del trabajo.

\section{Enfoque general}
\label{sec:enfoque-arquitectonico}

La arquitectura del sistema sigue un enfoque de ingeniería dirigida por modelos.
En lugar de definir manualmente cada bloque en JavaScript, el diseñador describe
el dominio mediante un metamodelo Ecore anotado o mediante una definición
textual \texttt{.m2b}. A partir de esa descripción, Model2Blockly deriva la
estructura y parte del comportamiento del editor. Esta organización aplica los
principios de MDE al diseño concreto de la solución: los modelos constituyen la
entrada principal del proceso y se utilizan para generar software
\cite{brambilla2017modelDriven}.

En ambas rutas se mantiene la separación habitual entre sintaxis abstracta y
sintaxis concreta de un DSL \cite{sal2024dsl}. Ecore o el lenguaje
\texttt{.m2b} describen los conceptos y relaciones del dominio, mientras que
Blockly proporciona la notación visual con la que el usuario final crea y edita
instancias. Por lo tanto, el lenguaje textual de Xtext no sustituye al editor
visual: es una forma compacta de especificar qué lenguaje de bloques debe
generarse.

El diseño se apoya en tres principios. El primero es separar la interpretación
del dominio de la generación del editor. El segundo es hacer que las dos rutas
de entrada produzcan un modelo \texttt{EditorSpec} conforme al mismo metamodelo
intermedio. Este modelo constituye el contrato arquitectónico compartido antes
de generar HTML y JavaScript. El tercero es derivar no solo la apariencia de los
bloques, sino también parte de su comportamiento: conexiones, validaciones,
referencias, exportación y metadatos de interfaz.

Estos principios responden al problema identificado en la introducción. Blockly
proporciona la infraestructura de interacción visual, pero exige que el
desarrollador describa cada bloque, campo, conexión y generador
\cite{blocklyDocs}. Model2Blockly desplaza esa descripción a un nivel de dominio:
el desarrollador define clases, atributos, contenciones, referencias,
anotaciones o construcciones textuales \texttt{.m2b}; la herramienta se encarga
de transformar esa información en los artefactos Blockly correspondientes.

\section{Flujo general de generación}
\label{sec:flujo-generacion}

El flujo completo de la solución se organiza en cuatro etapas:

\begin{enumerate}
  \item definición del dominio mediante un metamodelo Ecore con anotaciones o
  mediante un fichero textual \texttt{.m2b};
  \item interpretación de la entrada y transformación a un modelo intermedio
  común;
  \item generación de los artefactos Blockly a partir de ese modelo;
  \item ejecución del editor generado en el navegador.
\end{enumerate}

La Figura \ref{fig:arquitectura-general} muestra las dos rutas de entrada y la
cadena común que se ejecuta después de convertirlas al modelo intermedio.

\begin{figure}[H]
  \centering
  \begin{tikzpicture}[
      node distance=0.75cm and 2.2cm,
      box/.style={draw=black!55, rounded corners=2pt, minimum width=3.65cm,
        minimum height=0.92cm, align=center, font=\small, inner sep=5pt},
      inputbox/.style={box, fill=green!8},
      stagebox/.style={box, fill=blue!7},
      commonbox/.style={box, fill=orange!10},
      outputbox/.style={box, fill=orange!10},
      arrow/.style={-{Latex[length=2.5mm]}, thick, draw=black!70}
    ]
    \node[inputbox] (ecore) {Metamodelo Ecore\\anotado};
    \node[inputbox, right=of ecore] (dsl) {Definición textual\\\texttt{.m2b}};

    \node[stagebox, below=of ecore] (ecoreAdapter)
      {Transformación del\\modelo Ecore};
    \node[stagebox, below=of dsl] (xtextPath)
      {Análisis Xtext y transformación\\del modelo textual};

    \path let \p1=(ecoreAdapter), \p2=(xtextPath) in
      node[commonbox] (editorSpec) at ($ (\p1)!0.5!(\p2) +(0,-1.7cm) $)
      {Modelo intermedio común\\\texttt{EditorSpec}};
    \node[commonbox, below=of editorSpec] (generator)
      {Generación de artefactos\\Blockly};
    \node[outputbox, below=of generator] (outputs)
      {Editor Blockly ejecutable\\HTML/JavaScript};

    \draw[arrow] (ecore) -- (ecoreAdapter);
    \draw[arrow] (dsl) -- (xtextPath);
    \draw[arrow] (ecoreAdapter) -- (editorSpec);
    \draw[arrow] (xtextPath) -- (editorSpec);
    \draw[arrow] (editorSpec) -- (generator);
    \draw[arrow] (generator) -- (outputs);
  \end{tikzpicture}
  \caption{Arquitectura general de la generación de editores Blockly}
  \label{fig:arquitectura-general}
\end{figure}

La Figura \ref{fig:arquitectura-general} muestra dos recorridos diferentes en su
parte superior. En la ruta Ecore se interpreta la estructura del metamodelo y
sus anotaciones para construir el modelo intermedio. En la ruta textual, Xtext
analiza primero la definición \texttt{.m2b} y crea un modelo estructurado; una
segunda transformación lleva ese modelo a \texttt{EditorSpec}. Las dos rutas
terminan, por lo tanto, en modelos conformes al mismo metamodelo intermedio.

Por ejemplo, una \texttt{EClass Page} con un atributo
\texttt{title:EString} y una declaración textual \texttt{class Page} con
\texttt{attribute title:string} se representan de la misma forma después de la
conversión: un tipo de bloque \texttt{Page} con un campo de texto
\texttt{title}. Desde ese punto se aplica una única generación sobre
\texttt{EditorSpec}. Por ello, la solución necesita dos mecanismos de entrada,
pero comparte la representación intermedia y la generación de salida. Esta
convergencia es la decisión que estructura el resto de la arquitectura.

\section{Modelo intermedio}
\label{sec:modelo-intermedio}

Las dos rutas de entrada describen conceptos similares mediante representaciones
diferentes. Si la generación trabajara directamente con ambas, tendría que
mantener dos interpretaciones de clases, propiedades, relaciones y restricciones.
La incorporación de una nueva forma de entrada obligaría además a modificar la
generación. Para evitar esta dependencia, la arquitectura introduce un modelo
intermedio común entre las entradas y la salida.

Este modelo representa el editor que se desea obtener, no una copia completa del
artefacto de origen. Su contenido se organiza conceptualmente en tres grupos. El
primero describe la estructura general del editor, como las categorías de la
toolbox y su configuración. El segundo define los bloques, sus campos y sus
formas de conexión. El tercero conserva la información de comportamiento que
afecta a la edición, como referencias y reglas de validación. El metamodelo
intermedio establece las relaciones válidas entre estos conceptos y cada
transformación de entrada produce un modelo conforme a él.

La raíz de esta representación se denomina \texttt{EditorSpec}. Por ejemplo,
una clase de dominio \texttt{Page} se representa como un tipo de bloque y su
propiedad \texttt{title} como un campo editable. Una relación de composición se
representa como una zona para insertar bloques hijos, mientras que una referencia
a otro objeto se representa como un mecanismo de selección. El origen puede ser
Ecore o \texttt{.m2b}; después de la transformación, estos conceptos tienen la
misma forma para la generación.

% Cuando esté disponible el diagrama UML definitivo exportado desde PlantUML,
% insertarlo aquí con las clases, atributos, roles, multiplicidades y relaciones
% de contención del metamodelo intermedio.

La elección de esta representación intermedia permite que cada parte tenga una
responsabilidad definida. Las rutas de entrada interpretan el dominio y la
generación interpreta únicamente \texttt{EditorSpec}. Así se comparte una sola
lógica de salida y se pueden modificar o ampliar las entradas sin rediseñar el
proceso que produce el editor Blockly.

\section{Ruta basada en Ecore}
\label{sec:ruta-ecore}

La ruta Ecore está dirigida a dominios que ya se encuentran definidos dentro
del ecosistema EMF. Permite reutilizar un metamodelo existente como punto de
partida, en lugar de volver a describir manualmente sus conceptos para Blockly.
Ecore aporta la estructura abstracta del dominio mediante clases, atributos,
referencias, herencia y cardinalidades \cite{emfOverview}.

El mapeo básico es sistemático. Cada clase Ecore se convierte en un tipo de
bloque; cada atributo, en un campo visual; cada referencia de contención, en una
entrada estructural para bloques hijos; y cada referencia no contenida, en un
campo de referencia dinámica. La herencia se utiliza para mantener la
compatibilidad entre tipos de bloque y las cardinalidades permiten derivar
comprobaciones sobre elementos obligatorios o cantidades admitidas. Estas reglas
de correspondencia preservan en el editor las diferencias estructurales que son
relevantes para crear modelos válidos.

Ecore no describe por sí solo todas las decisiones de sintaxis concreta de un
editor Blockly, como el color, la etiqueta visible, la organización de la
toolbox o el tipo de control utilizado para editar un valor. Por este motivo, la
ruta admite anotaciones que complementan el metamodelo sin alterar su estructura
esencial. La separación entre estructura del dominio y metadatos de presentación
permite obtener una salida útil por defecto y personalizarla cuando el dominio
lo requiere.

\section{Ruta textual \texttt{.m2b} basada en Xtext}
\label{sec:ruta-xtext}

La ruta textual ofrece una alternativa de autoría cuando no se dispone de un
metamodelo Ecore previo o se prefiere una notación más compacta. Xtext resulta
adecuado porque permite definir una sintaxis textual y obtener a partir de ella
una representación estructurada del dominio \cite{xtextProject,bettini2013xtext}.

Esta ruta no reemplaza a Ecore, sino que responde a una necesidad distinta. El
DSL permite reunir en una misma descripción la estructura del dominio y las
decisiones básicas de presentación del editor. Xtext se encarga de interpretar
el texto y comprobar que respeta la sintaxis definida; posteriormente, esa
representación se transforma al mismo modelo intermedio empleado por la ruta
Ecore. La generación continúa siendo común y no depende de la forma elegida para
describir el dominio.

\section{Generación de editores Blockly}
\label{sec:generacion-editores}

Una vez construido el modelo intermedio, la generación produce los artefactos
necesarios para ejecutar el editor en un navegador. Esta etapa puede entenderse
como una transformación de modelo a texto: la especificación del editor se
convierte en definiciones de bloques, toolbox, configuración del workspace,
reglas de validación y páginas HTML y JavaScript.

La correspondencia se establece por responsabilidad: la organización del modelo
intermedio determina la toolbox, la definición de cada tipo determina la forma
del bloque y sus conexiones, y las reglas conservadas determinan las ayudas y
advertencias disponibles durante la edición. Así, la salida visual se deriva de
decisiones expresadas en el modelo y no de una colección de bloques programados
manualmente.

La dependencia exclusiva respecto al modelo intermedio evita duplicar la lógica
de salida para cada entrada y deja abierta la posibilidad de incorporar otros
formatos de origen o nuevos objetivos de generación. Para este trabajo se elige
como salida principal un editor HTML/JavaScript autocontenido, porque facilita
su ejecución y revisión sin exigir una infraestructura adicional.

\section{Validaciones, referencias y personalización}
\label{sec:validaciones-referencias-personalizacion}

Un editor de bloques no debe limitarse a mostrar piezas visuales; también debe
ayudar a construir modelos coherentes. Por ello, la arquitectura conserva en el
modelo intermedio las restricciones que pueden trasladarse al editor. Entre
ellas se encuentran las cardinalidades, la obligatoriedad de valores y algunas
relaciones sencillas entre elementos. La decisión de mantenerlas en el modelo
intermedio permite aplicar el mismo criterio de validación con independencia de
la ruta de entrada.

Estas comprobaciones no constituyen un motor general de restricciones. OCL
permite expresar invariantes sobre objetos, propiedades y colecciones
\cite{omgOcl}, mientras que Model2Blockly solo transforma patrones que pueden
reconocerse y ejecutarse de manera predecible en el editor generado. Esta
decisión mantiene explícito el alcance del sistema: las validaciones ayudan al
usuario mediante advertencias, pero no sustituyen a un lenguaje completo de
restricciones. La implementación concreta y la evaluación del subconjunto
soportado se presentan en los capítulos siguientes.

Las referencias no contenidas se tratan de forma distinta a las contenciones.
Una contención define la estructura jerárquica del modelo y por lo tanto se
representa como una entrada de sentencias. Una referencia no contenida apunta a
otro elemento existente en el workspace, por lo que se representa mediante un
desplegable dinámico. El generador calcula los candidatos compatibles teniendo
en cuenta el tipo objetivo y sus subtipos. Esta decisión mantiene una
diferencia semántica importante del metamodelo: no es lo mismo crear un hijo
estructural que seleccionar otro elemento ya presente en el modelo.

La personalización se incorpora como una capa controlada de metadatos. Colores,
etiquetas, categorías, tooltips, disposición en línea, widgets de interfaz y
plantillas de código no cambian la estructura esencial del dominio, pero sí
afectan a la experiencia del usuario final. Model2Blockly permite definir esta
información mediante anotaciones Ecore o mediante opciones del DSL textual. De
este modo, se evita que el usuario tenga que editar directamente JavaScript para
realizar cambios habituales de presentación.

\section{Trazabilidad con los objetivos}
\label{sec:trazabilidad-objetivos}

La arquitectura propuesta mantiene una relación directa con los objetivos
específicos definidos para el TFM. La Tabla
\ref{tab:trazabilidad-arquitectura} resume esa relación.

\begin{table}[h]
  \centering
  \begin{tabular}{|p{0.18\textwidth}|p{0.70\textwidth}|}
    \hline
    Objetivo & Respuesta arquitectónica \\
    \hline
    OE1 & El modelo intermedio identifica explícitamente los elementos que
    necesita un editor de bloques: tipos, campos, conexiones, categorías,
    referencias y validaciones. \\
    OE2 & La ruta Ecore permite partir de metamodelos y
    derivar bloques, campos, entradas y reglas. \\
    OE3 & La gramática Xtext proporciona la ruta textual \texttt{.m2b}
    para definir lenguajes de bloques de forma compacta. \\
    OE4 & El modelo intermedio común desacopla las entradas de los generadores
    y permite compartir la misma lógica de salida. \\
    OE5 & Los generadores producen editores HTML/JavaScript ejecutables en el
    navegador. \\
    OE6 & Las validaciones, referencias dinámicas, anotaciones de interfaz y
    plantillas de código se incorporan como capacidades genéricas. \\
    OE7 & La separación por transformación de entrada, modelo intermedio y
    generadores facilita probar cada parte y evaluar la convergencia de las dos
    rutas sobre AppMaker. \\
    \hline
  \end{tabular}
  \caption{Trazabilidad entre objetivos y arquitectura}
  \label{tab:trazabilidad-arquitectura}
\end{table}

En conjunto, la arquitectura define una cadena generativa clara: el diseñador
describe un dominio como metamodelo Ecore anotado o como fichero \texttt{.m2b},
el sistema transforma esa descripción en un modelo intermedio y el generador
produce un editor Blockly. La aportación principal no es un bloque concreto,
sino el mecanismo reutilizable que permite generar editores desde descripciones
de dominio de alto nivel.
